% !TeX root = ../../main.tex

\section{Soluciones de Big Data en la nube}

\subsection{Introducción}

Las plataformas de Big Data en la nube permiten almacenar, procesar y
analizar grandes cantidades de información sin que una organización tenga
que adquirir toda la infraestructura física. Sus recursos pueden aumentar o
disminuir de acuerdo con la demanda y normalmente se pagan según el uso.

Los servicios en la nube se pueden clasificar de la siguiente manera:

\begin{itemize}
    \item \textbf{IaaS:} proporciona servidores, redes y almacenamiento virtual.
    El cliente administra el sistema operativo y sus aplicaciones.
    
    \item \textbf{PaaS:} proporciona una plataforma administrada para desarrollar,
    procesar o analizar datos sin configurar toda la infraestructura.
    
    \item \textbf{SaaS:} ofrece una aplicación terminada a la que se accede
    normalmente mediante un navegador.
\end{itemize}

A continuación se presentan las principales soluciones de Big Data ofrecidas
por AWS, Google Cloud, Microsoft Azure, IBM, Cloudera y Databricks.

\subsection{Amazon Web Services}

\subsubsection*{Descripción y antecedentes}

Amazon Web Services, conocida como AWS, es la plataforma de servicios en la
nube de Amazon. Comenzó a ofrecer servicios públicos en 2006 y actualmente
cuenta con herramientas de almacenamiento, procesamiento, bases de datos,
analítica e inteligencia artificial.

Una de sus principales características es la elasticidad. Una empresa puede
aumentar sus recursos durante los periodos de mayor actividad y reducirlos
cuando ya no sean necesarios.

\subsubsection*{Tipos de servicio}

\begin{itemize}
    \item \textbf{IaaS:} Amazon EC2 proporciona servidores virtuales, mientras
    que Amazon S3 permite almacenar grandes cantidades de archivos y datos.
    
    \item \textbf{PaaS:} incluye servicios administrados como Amazon EMR,
    Amazon Redshift, AWS Glue, Amazon Athena y Amazon Kinesis.
    
    \item \textbf{SaaS:} Amazon QuickSight proporciona análisis y tableros
    visuales mediante una aplicación web.
\end{itemize}

\subsubsection*{Herramientas orientadas a Big Data}

\begin{itemize}
    \item \textbf{Amazon S3:} almacenamiento de objetos que puede utilizarse
    para construir un lago de datos.
    
    \item \textbf{Amazon EMR:} procesa grandes volúmenes de información con
    tecnologías como Apache Spark y Hadoop.
    
    \item \textbf{Amazon Redshift:} almacén de datos utilizado para realizar
    consultas y análisis empresariales.
    
    \item \textbf{AWS Glue:} prepara, transforma e integra información mediante
    procesos ETL.
    
    \item \textbf{Amazon Kinesis:} recibe y analiza datos en tiempo real.
    
    \item \textbf{Amazon Athena:} permite consultar archivos almacenados en S3
    utilizando instrucciones SQL.
\end{itemize}

\subsubsection*{Precios}

AWS utiliza principalmente el modelo de pago por uso. Cada servicio cobra de
manera diferente: S3 cobra por almacenamiento, EMR por los recursos de
procesamiento y Redshift por la capacidad utilizada. Como referencia, Athena
cobra aproximadamente \$5 dólares por cada TB de datos analizado. Los costos
pueden variar según la región y la configuración.

\subsubsection*{Ventajas y desventajas}

\textbf{Ventajas:}

\begin{itemize}
    \item Cuenta con una gran variedad de servicios.
    \item Permite aumentar o reducir los recursos rápidamente.
    \item Tiene disponibilidad en numerosas regiones.
    \item Sus herramientas se integran entre sí.
\end{itemize}

\textbf{Desventajas:}

\begin{itemize}
    \item Su catálogo puede resultar complicado para usuarios principiantes.
    \item Una mala configuración puede producir costos elevados.
    \item Migrar posteriormente a otra plataforma puede ser difícil.
\end{itemize}

\subsubsection*{Caso de éxito}

Netflix utiliza AWS para distribuir contenido a millones de usuarios y para
procesar información relacionada con reproducciones, dispositivos y
preferencias. El análisis de estos datos ayuda a mejorar las recomendaciones
y a conocer el comportamiento de los usuarios.

\subsubsection*{Video sugerido}

La página oficial del caso de Netflix contiene un video sobre su uso de AWS:

\url{https://aws.amazon.com/solutions/case-studies/netflix-case-study/}


\subsection{Google Cloud}

\subsubsection*{Descripción y antecedentes}

Google Cloud es la plataforma de servicios en la nube de Google. Sus primeros
servicios públicos aparecieron en 2008 y posteriormente se incorporaron
herramientas especializadas en análisis de datos, entre ellas BigQuery.

La plataforma aprovecha tecnologías desarrolladas por Google para procesar
información a gran escala. Se caracteriza por sus servicios administrados y
por reducir la necesidad de configurar servidores manualmente.

\subsubsection*{Tipos de servicio}

\begin{itemize}
    \item \textbf{IaaS:} Compute Engine ofrece máquinas virtuales y Cloud
    Storage proporciona almacenamiento de objetos.
    
    \item \textbf{PaaS:} BigQuery, Dataflow, Dataproc, Pub/Sub y Bigtable
    proporcionan procesamiento y análisis administrado.
    
    \item \textbf{SaaS:} Looker permite elaborar reportes, análisis y tableros
    empresariales.
\end{itemize}

\subsubsection*{Herramientas orientadas a Big Data}

\begin{itemize}
    \item \textbf{BigQuery:} almacén de datos sin servidor para analizar grandes
    conjuntos de información con SQL.
    
    \item \textbf{Dataflow:} procesa datos por lotes y en tiempo real.
    
    \item \textbf{Dataproc:} ejecuta tecnologías como Apache Spark y Hadoop.
    
    \item \textbf{Pub/Sub:} recibe mensajes y eventos producidos continuamente.
    
    \item \textbf{Bigtable:} base de datos NoSQL para aplicaciones que requieren
    manejar grandes cantidades de información.
    
    \item \textbf{Cloud Storage:} almacena archivos que pueden formar parte de
    un lago de datos.
\end{itemize}

\subsubsection*{Precios}

Google Cloud también utiliza pago por consumo. BigQuery puede cobrarse por los
datos procesados en cada consulta o mediante capacidad reservada. Como
referencia, ofrece el primer TiB procesado durante el mes sin costo y después
cobra aproximadamente \$6.25 dólares por TiB. El almacenamiento y otros
servicios se cobran por separado.

\subsubsection*{Ventajas y desventajas}

\textbf{Ventajas:}

\begin{itemize}
    \item BigQuery permite analizar datos sin administrar servidores.
    \item Tiene buenas herramientas para inteligencia artificial.
    \item Facilita el procesamiento de información en tiempo real.
    \item Puede integrarse con herramientas de código abierto.
\end{itemize}

\textbf{Desventajas:}

\begin{itemize}
    \item Las consultas mal diseñadas pueden aumentar el costo.
    \item Algunos servicios requieren conocimientos especializados.
    \item Tiene menos presencia empresarial que AWS o Azure en algunas regiones.
\end{itemize}

\subsubsection*{Caso de éxito}

Spotify utiliza Google Cloud para procesar información relacionada con
canciones, listas de reproducción y actividad de sus usuarios. La plataforma
le permite ejecutar miles de procesos de datos y generar recomendaciones
personalizadas.

\subsubsection*{Video sugerido}

Video oficial con una introducción y demostración de BigQuery:

\url{https://www.youtube.com/watch?v=BH_7_zVk5oM}


\subsection{Microsoft Azure}

\subsubsection*{Descripción y antecedentes}

Microsoft Azure es la plataforma de servicios en la nube de Microsoft. Fue
lanzada comercialmente en 2010 y se encuentra integrada con productos como
Windows Server, SQL Server, Microsoft 365 y Power BI.

Azure permite construir soluciones de análisis que combinan almacenamiento,
procesamiento, integración de información e inteligencia artificial.

\subsubsection*{Tipos de servicio}

\begin{itemize}
    \item \textbf{IaaS:} Azure Virtual Machines proporciona servidores y Azure
    Blob Storage ofrece almacenamiento de objetos.
    
    \item \textbf{PaaS:} Azure Synapse Analytics, Azure Data Factory, Event Hubs,
    Stream Analytics y HDInsight ofrecen servicios administrados.
    
    \item \textbf{SaaS:} Power BI permite crear reportes y tableros para apoyar
    la toma de decisiones.
\end{itemize}

\subsubsection*{Herramientas orientadas a Big Data}

\begin{itemize}
    \item \textbf{Azure Data Lake Storage:} almacena grandes cantidades de datos
    estructurados y no estructurados.
    
    \item \textbf{Azure Synapse Analytics:} integra almacenamiento de datos,
    consultas SQL y procesamiento con Apache Spark.
    
    \item \textbf{Azure Data Factory:} crea procesos para extraer, transformar
    y cargar información.
    
    \item \textbf{Azure Event Hubs:} recibe millones de eventos generados por
    aplicaciones y dispositivos.
    
    \item \textbf{Azure Stream Analytics:} analiza flujos de información en
    tiempo real.
    
    \item \textbf{Azure HDInsight:} proporciona clústeres administrados de
    Hadoop, Spark, Kafka y otras tecnologías.
\end{itemize}

\subsubsection*{Precios}

El costo depende del almacenamiento, las operaciones y la capacidad de
procesamiento. Synapse puede cobrar por la cantidad de datos consultados en su
modalidad sin servidor o por la capacidad de un grupo dedicado. Microsoft
también ofrece descuentos por reservar recursos durante uno o tres años.

\subsubsection*{Ventajas y desventajas}

\textbf{Ventajas:}

\begin{itemize}
    \item Se integra fácilmente con productos de Microsoft.
    \item Permite construir soluciones híbridas.
    \item Cuenta con herramientas de seguridad y administración empresarial.
    \item Power BI facilita la visualización de resultados.
\end{itemize}

\textbf{Desventajas:}

\begin{itemize}
    \item Su sistema de precios puede resultar difícil de calcular.
    \item La configuración inicial puede requerir personal especializado.
    \item Algunos servicios tienen funciones similares y pueden confundirse.
\end{itemize}

\subsubsection*{Caso de éxito}

Grupo Bimbo utiliza servicios como Azure Synapse Analytics y Power BI para
analizar grandes cantidades de transacciones. Esta información le ayuda a
mejorar sus rutas de distribución y a tomar decisiones comerciales con mayor
rapidez.

\subsubsection*{Video sugerido}

Video oficial con una explicación y demostración de Azure Synapse Analytics:

\url{https://www.youtube.com/watch?v=dvP0JwchjfI}


\subsection{IBM Cloud}

\subsubsection*{Descripción y antecedentes}

IBM Cloud reúne la experiencia de IBM en infraestructura, bases de datos,
analítica e inteligencia artificial. Su oferta fue fortalecida con servicios
como SoftLayer, Bluemix y, posteriormente, la familia watsonx.

Su enfoque principal son las empresas que necesitan trabajar en ambientes
híbridos, es decir, combinar sistemas propios con recursos de nube pública.

\subsubsection*{Tipos de servicio}

\begin{itemize}
    \item \textbf{IaaS:} IBM Virtual Servers y Cloud Object Storage proporcionan
    procesamiento y almacenamiento.
    
    \item \textbf{PaaS:} watsonx.data, Db2 Warehouse, Analytics Engine y Event
    Streams permiten procesar y analizar información.
    
    \item \textbf{SaaS:} IBM Cognos Analytics proporciona reportes,
    visualizaciones y análisis empresarial.
\end{itemize}

\subsubsection*{Herramientas orientadas a Big Data}

\begin{itemize}
    \item \textbf{IBM Cloud Object Storage:} almacena datos no estructurados.
    
    \item \textbf{IBM watsonx.data:} plataforma abierta de tipo lakehouse para
    consultar y gobernar datos empresariales.
    
    \item \textbf{IBM Db2 Warehouse:} almacén de datos para análisis mediante SQL.
    
    \item \textbf{IBM Analytics Engine:} proporciona procesamiento con Apache
    Spark.
    
    \item \textbf{IBM Event Streams:} procesa eventos en tiempo real con
    tecnología basada en Apache Kafka.
    
    \item \textbf{IBM Cognos Analytics:} crea reportes y tableros empresariales.
\end{itemize}

\subsubsection*{Precios}

IBM utiliza precios por consumo y planes empresariales. En watsonx.data los
recursos se miden mediante unidades de recurso o RU. Como referencia, algunos
planes utilizan un precio aproximado de \$1 dólar por RU. El almacenamiento,
la transferencia y otros componentes pueden cobrarse por separado.

\subsubsection*{Ventajas y desventajas}

\textbf{Ventajas:}

\begin{itemize}
    \item Tiene experiencia en soluciones para grandes empresas.
    \item Su enfoque híbrido permite trabajar con datos locales y en la nube.
    \item Ofrece herramientas de gobierno y seguridad de la información.
    \item Utiliza tecnologías y formatos abiertos.
\end{itemize}

\textbf{Desventajas:}

\begin{itemize}
    \item Su catálogo es menor que el de AWS, Azure o Google Cloud.
    \item Algunas soluciones están orientadas principalmente a empresas grandes.
    \item Puede requerir asesoría especializada para su implementación.
\end{itemize}

\subsubsection*{Caso de éxito}

IBM utiliza sus servicios de datos e inteligencia artificial en eventos
deportivos como Wimbledon. La información de partidos y jugadores se procesa
rápidamente para producir estadísticas y contenido que mejora la experiencia
de los aficionados.

\subsubsection*{Video sugerido}

Biblioteca oficial de demostraciones de IBM watsonx.data:

\url{https://www.ibm.com/products/watsonx-data/demo-library}


\subsection{Cloudera y Hortonworks}

\subsubsection*{Descripción y antecedentes}

Cloudera y Hortonworks fueron empresas especializadas en plataformas basadas
en Apache Hadoop. Cloudera fue fundada en 2008 y Hortonworks en 2011. Ambas
empresas completaron su fusión en enero de 2019.

Por esta razón, Hortonworks ya no funciona como una plataforma independiente.
Sus tecnologías se integraron en la actual Cloudera Data Platform, conocida
como CDP.

\subsubsection*{Tipos de servicio}

\begin{itemize}
    \item \textbf{IaaS:} Cloudera no ofrece infraestructura general propia. La
    plataforma utiliza los recursos de AWS, Azure, Google Cloud o servidores
    locales.
    
    \item \textbf{PaaS:} CDP proporciona servicios administrados de ingeniería,
    almacenamiento, procesamiento, gobierno y aprendizaje automático.
    
    \item \textbf{SaaS:} algunas funciones administradas de análisis y
    visualización se ofrecen mediante acceso web.
\end{itemize}

\subsubsection*{Herramientas orientadas a Big Data}

\begin{itemize}
    \item \textbf{Cloudera Data Engineering:} ejecuta procesos de transformación
    mediante Apache Spark.
    
    \item \textbf{Cloudera Data Warehouse:} permite realizar consultas con
    tecnologías como Hive e Impala.
    
    \item \textbf{Cloudera DataFlow:} mueve y procesa información con Apache
    NiFi y Kafka.
    
    \item \textbf{Cloudera Data Hub:} permite crear clústeres para diferentes
    cargas de trabajo.
    
    \item \textbf{Cloudera Machine Learning:} ayuda a desarrollar y publicar
    modelos de aprendizaje automático.
    
    \item \textbf{Cloudera SDX:} administra seguridad, catálogo, permisos y
    gobierno de los datos.
\end{itemize}

\subsubsection*{Precios}

En la nube pública, Cloudera cobra mediante unidades de cómputo llamadas CCU.
Dependiendo del servicio, algunos precios de referencia se encuentran entre
\$0.04 y \$0.20 dólares por CCU por hora. A este precio se debe sumar la
infraestructura utilizada en AWS, Azure o Google Cloud. Las instalaciones
locales normalmente se contratan mediante una suscripción empresarial.

\subsubsection*{Ventajas y desventajas}

\textbf{Ventajas:}

\begin{itemize}
    \item Puede trabajar en la nube y en centros de datos propios.
    \item Mantiene compatibilidad con muchas tecnologías de código abierto.
    \item Proporciona gobierno y seguridad centralizados.
    \item Facilita la migración de sistemas Hadoop existentes.
\end{itemize}

\textbf{Desventajas:}

\begin{itemize}
    \item Su administración puede ser compleja.
    \item Requiere conocimientos de Hadoop, Spark y otras tecnologías.
    \item El costo incluye la licencia y la infraestructura del proveedor cloud.
\end{itemize}

\subsubsection*{Caso de éxito}

Podo, una empresa española del sector energético, utiliza Cloudera y Google
Cloud para analizar información de consumo eléctrico. El procesamiento de
estos datos le permite generar pronósticos y mejorar la administración de la
energía.

\subsubsection*{Video sugerido}

Demostración oficial de ingeniería de datos con Cloudera:

\url{https://video.cloudera.com/detail/video/6324827111112/}


\subsection{Databricks}

\subsubsection*{Descripción y antecedentes}

Databricks fue fundada en 2013 por varios de los creadores de Apache Spark.
Su propuesta se basa en el concepto de \textit{lakehouse}, el cual combina las
capacidades de un lago de datos y un almacén de datos.

La plataforma permite que analistas, ingenieros de datos y científicos de
datos trabajen con la misma información y dentro de un entorno compartido.

\subsubsection*{Tipos de servicio}

\begin{itemize}
    \item \textbf{IaaS:} Databricks utiliza la infraestructura de AWS, Azure o
    Google Cloud, pero no ofrece una nube de infraestructura propia.
    
    \item \textbf{PaaS:} ofrece entornos administrados para ingeniería de datos,
    analítica, inteligencia artificial y aprendizaje automático.
    
    \item \textbf{SaaS:} Databricks SQL, los tableros y las herramientas de
    colaboración se utilizan mediante una interfaz web.
\end{itemize}

\subsubsection*{Herramientas orientadas a Big Data}

\begin{itemize}
    \item \textbf{Apache Spark:} procesa grandes cantidades de datos de manera
    distribuida.
    
    \item \textbf{Delta Lake:} agrega confiabilidad, historial y control de
    transacciones a un lago de datos.
    
    \item \textbf{Databricks SQL:} permite consultar datos y crear tableros.
    
    \item \textbf{Unity Catalog:} administra permisos, metadatos y gobierno de
    la información.
    
    \item \textbf{MLflow:} ayuda a administrar el ciclo de vida de los modelos
    de aprendizaje automático.
    
    \item \textbf{Notebooks colaborativos:} permiten trabajar con Python, SQL,
    R y Scala.
\end{itemize}

\subsubsection*{Precios}

El consumo se mide mediante unidades llamadas DBU. El costo total depende del
tipo de carga de trabajo, las DBU utilizadas y la infraestructura contratada
en AWS, Azure o Google Cloud. Existen modalidades de pago por uso, contratos
con consumo comprometido y opciones sin servidor.

\subsubsection*{Ventajas y desventajas}

\textbf{Ventajas:}

\begin{itemize}
    \item Integra ingeniería de datos, SQL e inteligencia artificial.
    \item Facilita la colaboración entre diferentes equipos.
    \item Puede utilizarse en las tres principales plataformas cloud.
    \item Está construido alrededor de Apache Spark y formatos abiertos.
\end{itemize}

\textbf{Desventajas:}

\begin{itemize}
    \item El costo puede aumentar si los clústeres permanecen activos.
    \item Requiere conocimientos de programación y procesamiento distribuido.
    \item Depende de la infraestructura de otro proveedor cloud.
\end{itemize}

\subsubsection*{Caso de éxito}

Shell utiliza Databricks como parte de su estrategia de análisis e
inteligencia artificial. La plataforma permite que sus equipos preparen datos,
entrenen modelos y colaboren en proyectos relacionados con energía,
mantenimiento y optimización de operaciones.

\subsubsection*{Video sugerido}

Demostración oficial de aprendizaje automático con Databricks:

\url{https://www.youtube.com/watch?v=1TPu-uOLglI}


\subsection{Comparación general}

\begin{itemize}
    \item \textbf{AWS} destaca por la cantidad y variedad de sus servicios.
    
    \item \textbf{Google Cloud} resulta conveniente para análisis sin servidor,
    principalmente mediante BigQuery.
    
    \item \textbf{Microsoft Azure} es útil para organizaciones que ya emplean
    productos de Microsoft.
    
    \item \textbf{IBM Cloud} se orienta a soluciones empresariales híbridas y
    al gobierno de la información.
    
    \item \textbf{Cloudera} es apropiada para ambientes híbridos y organizaciones
    que utilizan Hadoop y otras tecnologías abiertas.
    
    \item \textbf{Databricks} ofrece una plataforma unificada para ingeniería de
    datos, análisis e inteligencia artificial.
\end{itemize}

\subsection{Conclusión}

No existe una plataforma que sea la mejor para todos los casos. La elección
depende del presupuesto, la experiencia del equipo, las herramientas que ya
utiliza la organización y la cantidad de información que necesita procesar.

Las plataformas estudiadas permiten almacenar y analizar grandes volúmenes de
datos, pero tienen diferencias importantes. AWS y Azure ofrecen catálogos muy
amplios; Google Cloud facilita el análisis sin servidor; IBM y Cloudera se
enfocan en soluciones híbridas; mientras que Databricks integra ingeniería de
datos e inteligencia artificial en un mismo entorno. Antes de seleccionar una
solución es necesario realizar una prueba y calcular su costo estimado.